\section{Frontier RL Flywheel 与能力边界}

有了真实环境，系统可以生成任务、尝试行动、验证结果，再把成功轨迹用于更新模型。讲者把这一过程称为 recursive self-improvement，但这个词很容易被误读成模型完全脱离人类、无限自主改进。更准确的系统解释是：模型帮助扩大实验与数据生产，人类和基础设施仍负责选择环境、定义验证、控制部署、处理异常并决定哪些改进值得进入下一轮。

\begin{figure}[H]
\centering
\includegraphics[width=0.94\textwidth]{images/00-35-15-frontier-rl-flywheel.jpg}
\caption{Frontier RL flywheel：从假设与环境开始，经过尝试、验证和学习，再生成更强的下一轮系统。}
\figsource{Stanford Online 官方视频 00:35:15，`images/00-35-15-frontier-rl-flywheel.jpg`。}
\end{figure}

\begin{knowledgebox}{读图：递归改进需要外部锚点}
左侧是任务、数据或科学假设的生成，中间是模型在环境中执行并得到可验证结果，右侧是把有效经验写回策略、工具或知识库。闭环中的外部锚点包括单元测试、物理实验、用户验收、安全规则与成本上限。若没有可靠锚点，系统可能只是在优化代理指标，形成看似快速、实际偏离目标的自我强化。
\end{knowledgebox}

\subsection{边界一：经验上还能扩多久}

经验边界关心的是可观察曲线：增加 rollout、环境数量、任务难度和训练算力后，成功率是否继续提升，边际收益何时下降，训练是否出现模式坍缩或 reward hacking。这个问题可以通过实验逐步回答。不同领域的答案也会不同，代码测试、数学答案、材料实验和开放式写作拥有完全不同的验证成本。

\begin{figure}[H]
\centering
\includegraphics[width=0.9\textwidth]{images/00-37-15-limits-to-rl.jpg}
\caption{讲者提出 RL 的两类边界：经验上的平台期，以及关于目标与人类价值的哲学边界。}
\figsource{Stanford Online 官方视频 00:37:15，`images/00-37-15-limits-to-rl.jpg`。}
\end{figure}

\begin{warningbox}{读图：两类问题不能用同一条曲线回答}
``更多算力后准确率是否继续上升'' 是经验问题，可以设计对照实验；``应该优化什么、谁有权定义成功、哪些能力不应部署'' 是规范与治理问题。后者不会因为 benchmark 继续上涨而自动消失。把两者混在一起，会让技术曲线替代价值选择，也会让价值争论忽略真实工程证据。
\end{warningbox}

\subsection{边界二：reward 代表谁的目标}

哲学边界最终会变成系统接口。谁选择训练任务，谁编写 verifier，谁能修改 reward，失败由谁承担，用户是否知道系统在优化什么，这些都不是训练结束后的附加问题。RL 的优势是把反馈纳入学习；它的风险也来自同一点：一旦反馈口径偏窄、可被操纵或代表性不足，优化会把偏差放大成稳定行为。

\begin{table}[H]
\centering
\begin{tabular}{L{0.2\textwidth}L{0.31\textwidth}L{0.39\textwidth}}
\toprule
边界类型 & 可观察信号 & 应对方式 \\
\midrule
算力边界 & rollout 吞吐、训练时长、能源与成本 & 系统优化、采样效率、资源预算 \\
环境边界 & 任务覆盖不足、模拟与现实偏差 & 扩展真实环境、分层验证、人工审计 \\
验证边界 & reward hacking、测试泄漏、代理指标失真 & 多重 verifier、隐藏测试、失败案例库 \\
组织边界 & 责任模糊、上线压力覆盖安全信号 & 明确决策权、事故复盘、停止条件 \\
价值边界 & 不同群体对成功定义冲突 & 参与式设计、制度约束、可申诉机制 \\
\bottomrule
\end{tabular}
\caption{把 ``RL 的极限'' 拆成可测量、可治理的不同问题。}
\end{table}

\begin{knowledgebox}{课堂提示：先亲手跑一次闭环}
讲者鼓励学生通过 problem set 理解 RL，不要只从宏观叙事判断。亲手定义环境、写 reward、观察策略钻漏洞，会迅速暴露“目标看起来清楚”和“目标可被可靠计算”之间的距离。这种实践也是 preparedness：知道什么时候该相信曲线，什么时候该检查数据和 verifier。
\end{knowledgebox}

\subsection{本章小结}

Frontier RL flywheel 可以扩大可验证经验的生产，但每一轮都依赖环境、verifier、人类选择和基础设施。经验边界可以用曲线和实验探索，哲学与治理边界需要明确责任和公共讨论。下面把视角从单个训练闭环拉到宏观市场，观察资本开支、软件采用和算力价格是否支持“大转型”叙事。

